% SOURCE_AUTHORITY: sga5_fr_workpass.tex, lines 4872--4907 (LF-normalized unit).
% SOURCE_SHA256: 791F4EFFC5E02832D5D77ED03518C8156D6F07E4C8238B03545DB93D883FBB28
Fica, portanto, demonstrada a). Demonstremos b). Seja $\bar\eta$ um ponto genérico geométrico de $Y$. Escolhamos flechas de especialização
\[
a:\bar\eta\to Y(\bar z),\qquad b:\bar\eta\to Y(\bar s),
\]
onde $Y(\bar z)$ (resp. $Y(\bar s)$) é a localização estrita de $Y$ em $\bar z$ (resp. $\bar s$). Ponhamos
\[
M=m^*\R j_*(i'_!L).
\]
O quadrado comutativo
\[
\begin{tikzcd}
Y & Y(\bar s)\arrow[l]\\
Y(\bar z)\arrow[u] & \bar\eta \arrow[l,"a"']\arrow[u,"b"']
\end{tikzcd}
\]
induz um quadrado comutativo
\[
\begin{tikzcd}[column sep=large,row sep=large]
\R\Gamma(Y,M) \arrow[r,"4.3.1"] \arrow[d] & M_{\bar s} \arrow[d,"b^*"]\\
M_{\bar z} \arrow[r,"a^*"'] & M_{\bar\eta}.
\end{tikzcd}
\]
Como $M|Y-z$ tem cohomologia localmente constante, $b^*$ é um isomorfismo. Por outro lado, como $M=i_{1!}i_1^*M$, tem-se $M_{\bar z}=0$, portanto a flecha 4.3.1 é nula. Resta demonstrar c). Tem-se $\widetilde C=C\times_Z\widetilde Z$, onde $\widetilde Z$ designa, como antes, a localização estrita de $Z$ em $\bar z$. Ponhamos $\widetilde Y=Y\times_Z\widetilde Z$. Tem-se então um diagrama comutativo de flechas de restrição
\[
\begin{tikzcd}[column sep=large,row sep=large]
\R\Gamma(Z-Y,i'_!L)=\R\Gamma(Z,\R j_*(i'_!L))
 \arrow[r,"4.3.2"] \arrow[d]
&
\R\Gamma(\widetilde C-\bar s,i'_!(L)|\widetilde C-\bar s)\\
\R\Gamma(Y,m^*\R j_*(i'_!L)) \arrow[r,"4.3.1"']
&
(\R j_*(i'_!L))_{\bar s}
  =\R\Gamma(\widetilde Z-\widetilde Y,(i'_!L)|\widetilde Z-\widetilde Y) \arrow[u]
\end{tikzcd}
\]
A nulidade de 4.3.2 deduz-se, portanto, de b), o que conclui a demonstração de 4.3.
